\documentclass[]{beamer}
%\usepackage[utopia]{mathdesign}
%\usepackage[no-math]{fontspec}
%\setmainfont{Liberation Serif}

\usepackage{minted}
% \usepackage{redhat}
\usepackage{hyperref}
\usepackage{ccicons}
\usepackage{ulem}

\usepackage{tikz}
\usetikzlibrary{arrows,shapes,snakes,automata,backgrounds,petri}

%gets rid of bottom navigation bars
\setbeamertemplate{footline}[page number]{}

%gets rid of navigation symbols
\setbeamertemplate{navigation symbols}{}

\newcommand{\done}{\textcolor{teal}{\checkmark}}
\newcommand{\pull}[3]{\href{https://github.com/systemd/#1/pulls/#2}{#3 (\##2)}}
\newcommand{\pulldone}[3]{\pull{#1}{#2}{#3} \done}
\newcommand{\commit}[3]{\href{https://github.com/systemd/#1/commit/#2}{#3 (#2)}}
\newcommand{\commitdone}[3]{\commit{#1}{#2}{#3} \done}

\title{Securing Python services using systemd}
\author{Zbigniew Jędrzejewski-Szmek}
\institute{%
  \includegraphics[width=0.4\textwidth]{images/Logo-redhat-color-375.png}\\
  \medskip
  \textit{zbyszek@in.waw.pl}\\
  \medskip
  \ccbysa
}
\date{\tiny PyConPL, 5.11.2022} % fifth day of the month of falling leaves, the third year of the pandemic

\begin{document}
\begin{frame}
\titlepage % Print the title page as the first slide
\end{frame}

\begin{frame}
  \frametitle{Agenda}

  \begin{itemize}
  \item systemd — running services
  \item python — example service
  \item systemd — unit protections
  \item hardening the example service
  \item hardening the example service more
  \item passing zecure sekretz to the service
  \item (pystemd — starting transient services)
  \end{itemize}
\end{frame}

\begin{frame}
  \frametitle{systemd units}

  traditional approach:\\ \texttt{\color{brown} /usr/lib/systemd/system/foobar.service}

  \quad

  alternative approach:\\ \texttt{\color{teal} systemd-run -u foobar.service …}
\end{frame}

\begin{frame}
  \frametitle{\texttt{tonietak.py}}

  \hfill\href{https://github.com/keszybz/securing-pyservices-using-systemd/blob/main/service/tonietak.py}{\includegraphics[width=0.5\textwidth]{images/tonietak.png}}\hfill{}

  \hfill\href{https://github.com/keszybz/securing-pyservices-using-systemd/blob/main/service/tonietak.py}{github.com/keszybz/…/tonietak.py}\hfill{}
\end{frame}

\begin{frame}
  \frametitle{Unit protections}

  \begin{itemize}
  \item User=
  \item DynamicUser=yes

  \item StateDirectory=yes
  \item WorkingDirectory=yes

  \end{itemize}
\end{frame}

\begin{frame}
  \frametitle{Unit protections}

  \begin{itemize}
  \item ProtectHome=yes
  \item ProtectSystem=yes
  \item NoExecPaths=
  \end{itemize}
\end{frame}

\begin{frame}
  \frametitle{Unit protections — more}

  \begin{itemize}
  \item ProtectTmp=yes
  \item PrivateNetworking=yes
  \item SystemCallFilter=\\
        \texttt{systemd-analyze syscall-filter}
    
  \end{itemize}
\end{frame}

\begin{frame}
  \frametitle{Encrypted Secrets}

  \begin{itemize}
  \item \texttt{systemd-creds encrypt}
  \item \texttt{LoadCredentialEncrypted=}

  \end{itemize}
\end{frame}

\begin{frame}[fragile]
  \frametitle{Finding ideas for more protection directives}

  \begin{minted}{bash}
systemd-analyze security run-u4020.service
  \end{minted}
\end{frame}

\begin{frame}
  \frametitle{Putting the service in production}

  Running from the command-line is not a deployment strategy.

  How to make the unit permanent?
\end{frame}

\begin{frame}[fragile]
  \frametitle{Strategy I: save file}

  \begin{minted}{bash}
systemctl cat run-u3563.service

src=$(systemctl show -P FragmentPath run-u3563.service)
sudo cp "$src" \
    /etc/systemd/system/flask-app-tonietak.service
\end{minted}
\end{frame}

\begin{frame}[fragile]
  \frametitle{Strategy II: \texttt{pystemd}}

  \begin{minted}{python}
import pystemd.run

cmd = 'flask --app tonietak run'.split()

properties = dict(
  DynamicUser=True,
  StateDirectory='files',
  WorkingDirectory='/var/lib',
  ProtectHome='yes',
  ProtectSystem='yes',
  # LoadCredentialEncrypted='password.bin',
)

pystemd.run(cmd, extra=properties)
  \end{minted}
\end{frame}

\begin{frame}[fragile]
  \frametitle{Acknowledgments and links}

  Kushal Das — ``Securing Services Using systemd''\\
  \url{https://www.youtube.com/watch?v=UUW8RO4hkg0},
  \href{https://github.com/kushaldas/verybad/blob/main/src/main.rs}{verybad.rs}

  \vfill

  \url{https://flask.palletsprojects.com/en/2.2.x/}{flask}\\
  (\textcolor{gray}{with apologies to all flask developers ;)})

  \vfill

  \href{https://www.freedesktop.org/software/systemd/man/systemd.exec.html}{systemd.exec(5)},
  \href{https://www.freedesktop.org/software/systemd/man/systemd.directives.html}{systemd.directives(5)}

  \vfill

  \href{https://pypi.org/project/pystemd/}{pystemd},
  \href{https://github.com/facebookincubator/pystemd/blob/master/_docs/pystemd.run.md}{pystemd.run}

  \vfill

  \textcolor{teal}{\href{https://raw.githubusercontent.com/keszybz/securing-pyservices-using-systemd/main/systemd-python-services.pdf}{systemd-python-services.pdf}} (these slides),
  \textcolor{teal}{\href{https://github.com/keszybz/securing-pyservices-using-systemd/blob/main/service/tonietak.py}{tonietak.py}}
\end{frame}


\end{document}
